\documentclass[12pt]{article}

\usepackage{amsmath,amssymb,amsfonts,amsthm}
\usepackage{hyperref}
\usepackage{geometry}
\usepackage{authblk}
\geometry{margin=1in}

\newtheorem{theorem}{Theorem}[section]
\newtheorem{proposition}[theorem]{Proposition}
\newtheorem{lemma}[theorem]{Lemma}
\newtheorem{corollary}[theorem]{Corollary}
\theoremstyle{definition}
\newtheorem{definition}[theorem]{Definition}
\theoremstyle{remark}
\newtheorem{remark}[theorem]{Remark}

\newcommand{\R}{\mathbb{R}}
\newcommand{\Z}{\mathbb{Z}}
\newcommand{\Q}{\mathbb{Q}}
\newcommand{\SL}{\mathrm{SL}}
\newcommand{\SO}{\mathrm{SO}}
\newcommand{\GL}{\mathrm{GL}}
\newcommand{\Sym}{\mathrm{Sym}}
\newcommand{\Symp}{\mathrm{Sym}^{+}}
\newcommand{\Symz}{\mathrm{Sym}_{0}}
\newcommand{\Diagz}{\mathrm{Diag}_{0}}
\newcommand{\calM}{\mathcal{S}}
\newcommand{\tr}{\operatorname{tr}}
\newcommand{\diag}{\operatorname{diag}}
\newcommand{\ad}{\operatorname{ad}}

\title{Scale-Free Quadratic Forms, Symmetric Space Geometry,\\
and Arithmetic Logarithmic Lattices}
\author[1]{Marvin L.\ Gentry}
\affil[1]{Independent Researcher\\
  \texttt{drmlgentry@protonmail.com}\\
  ORCID: \href{https://orcid.org/0009-0006-4550-2663}{0009-0006-4550-2663}}
\date{}

\begin{document}
\maketitle

\begin{abstract}
The space of positive definite quadratic forms modulo overall scaling
arises naturally whenever only ratios of quantities carry intrinsic
meaning.
We prove that this scale-free moduli space $\calM_n$ is canonically
diffeomorphic to the Riemannian symmetric space $\SL(n,\R)/\SO(n)$,
a complete simply connected manifold of nonpositive sectional curvature.
Its tangent space at the identity consists of traceless symmetric
matrices of dimension $\frac{n(n+1)}{2}-1$.
We clarify the distinction between quadratic deformations
and arbitrary linear endomorphisms:
the latter yield $n^2-1$, while quadratic data admit only symmetric
deformations.
We develop the Riemannian geometry of $\calM_n$:
Cartan decomposition, sectional curvature,
geodesics, Weyl chambers, and maximal flats.
Arithmetic logarithmic unit lattices of totally real
number fields embed discretely into maximal totally geodesic flats,
and the Dirichlet regulator is identified with Euclidean covolume.
A flat rigidity theorem shows that maximal flats are uniquely
determined by their tangent Cartan subalgebras.
Explicit examples illustrate the theory.
\end{abstract}

%--------------------------------------------------
\section{Introduction}
\label{sec:intro}
%--------------------------------------------------

Many geometric and arithmetic quantities are inherently multiplicative:
mass ratios, eigenvalue ratios, regulator volumes, and Gram matrix
determinants all carry meaning only through their ratios, not their
absolute values.
Passing to logarithmic coordinates converts multiplicative structure
to additive structure, but introduces a redundancy: an overall additive
shift corresponds to simultaneous rescaling of all quantities.
Factoring out this redundancy yields a \emph{scale-free} moduli space
that captures only the intrinsic shape of the data.

For data organized as a positive definite quadratic form of rank $n$,
the appropriate scale-free space is the quotient $\calM_n$ of positive
definite symmetric matrices by positive scalar multiples.
A natural question is: what is the correct geometric structure of
$\calM_n$, and in particular, how many independent scale-free
deformation directions does it admit?

A common answer in the literature is $n^2 - 1$, arising from the
dimension of the space of traceless linear endomorphisms
$\mathrm{End}_{\tr=0}(V)$.
This count applies to general linear deformations of a vector
space, but it does not apply in the quadratic setting.
Quadratic forms are symmetric bilinear objects, and their infinitesimal
deformations must preserve symmetry.
The correct space of scale-free deformation directions is therefore
the space of traceless \emph{symmetric} matrices $\Symz(n,\R)$, of
dimension
\[
  \frac{n(n+1)}{2} - 1.
\]
For $n = 3, 4, 5$ this gives $5, 9, 14$ respectively, not $8, 15, 24$.

The purpose of this paper is threefold.
First, we prove rigorously that $\calM_n \cong \SL(n,\R)/\SO(n)$
and that its tangent space is $\Symz(n,\R)$.
Second, we develop the complete Riemannian geometry of this symmetric
space: Cartan decomposition, curvature, geodesics, Weyl chambers, and
maximal flats.
Third, we embed arithmetic logarithmic unit lattices of totally real
number fields as discrete subsets of maximal flats, identifying the
Dirichlet regulator as a Euclidean covolume.
This provides a geometric interpretation of classical arithmetic
invariants and a foundation for further arithmetic-geometric
investigations.

The paper is organized as follows.
Sections~\ref{sec:forms}--\ref{sec:scale} establish the symmetric
space identification.
Section~\ref{sec:tangent} computes the tangent space and corrects
the dimensional count.
Sections~\ref{sec:cartan}--\ref{sec:weyl} develop the Riemannian
geometry.
Sections~\ref{sec:lattice}--\ref{sec:flats} embed arithmetic unit
lattices and prove the flat rigidity theorem.
Section~\ref{sec:examples} works out explicit examples.
Section~\ref{sec:conclusion} summarizes the results.

The geometric construction developed here provides the theoretical
foundation for the companion results in~\cite{Gentry:CKM,Gentry:PMNS},
where the CKM and PMNS mixing matrices are derived explicitly from the
$K$- and $N$-factors of the Iwasawa decomposition of compact hyperbolic
3-manifold holonomy~\cite{Gentry:CKM,Gentry:PMNS}.

%--------------------------------------------------
\section{Positive Definite Quadratic Forms}
\label{sec:forms}
%--------------------------------------------------

Let $V$ be a real vector space of dimension $n$.
A \emph{positive definite quadratic form} on $V$ is a symmetric
bilinear map $Q : V \times V \to \R$ satisfying $Q(v,v) > 0$ for
all nonzero $v \in V$.
In a chosen basis, $Q$ is represented by a matrix in
\[
  \Symp(n,\R)
  = \{G \in M_n(\R) : G = G^T,\; v^T G v > 0\ \forall v \neq 0\},
\]
the open cone of positive definite symmetric matrices.

The group $\GL(n,\R)$ acts transitively on $\Symp(n,\R)$ by
\[
  g \cdot G = g G g^T.
\]

\begin{proposition}
\label{prop:homogeneous}
$\Symp(n,\R)$ is a smooth homogeneous manifold diffeomorphic to
$\GL(n,\R)/\mathrm{O}(n)$.
\end{proposition}

\begin{proof}
Every $G \in \Symp(n,\R)$ admits a unique Cholesky decomposition
$G = \ell\ell^T$ with $\ell$ lower triangular and positive diagonal
entries.
The map $g \mapsto g g^T$ is surjective from $\GL(n,\R)$ onto
$\Symp(n,\R)$.
The stabilizer of the identity $I$ under the action
$g \cdot G = gGg^T$ is $\{g : gg^T = I\} = \mathrm{O}(n)$.
\end{proof}

The tangent space of $\Symp(n,\R)$ at $I$ is all of $\Sym(n,\R)$,
the space of symmetric matrices, of dimension $\frac{n(n+1)}{2}$.

%--------------------------------------------------
\section{Scale Quotient and Symmetric Space Identification}
\label{sec:scale}
%--------------------------------------------------

Overall scaling acts on $\Symp(n,\R)$ by $G \mapsto \lambda G$ for
$\lambda > 0$.
Two forms that differ by a positive scalar factor represent the same
shape.

\begin{definition}[Scale-free quadratic moduli space]
The \emph{shape space} of rank $n$ is
\[
  \calM_n = \Symp(n,\R)/\R_{>0}.
\]
\end{definition}

\begin{proposition}
\label{prop:det1slice}
The map $G \mapsto (\det G)^{-1/n} G$ identifies $\calM_n$
diffeomorphically with the determinant-one slice
\[
  \Symp_1(n,\R)
  = \{G \in \Symp(n,\R) : \det G = 1\}.
\]
Each equivalence class in $\calM_n$ contains a unique representative
in $\Symp_1(n,\R)$.
\end{proposition}

\begin{proof}
For any $G \in \Symp(n,\R)$, set $\lambda = (\det G)^{1/n} > 0$.
Then $\det(\lambda^{-1} G) = \lambda^{-n} \det G = 1$, so
$\lambda^{-1} G \in \Symp_1(n,\R)$.
Two matrices $G, G'$ lie in the same class iff $G' = \mu G$ for some
$\mu > 0$; their normalized representatives coincide iff
$(\det G')^{-1/n} G' = (\det G)^{-1/n} G$, which holds iff $G' = \mu G$.
\end{proof}

\begin{theorem}
\label{thm:symspace}
There is a canonical $\SL(n,\R)$-equivariant diffeomorphism
\[
  \calM_n \;\cong\; \SL(n,\R)/\SO(n).
\]
\end{theorem}

\begin{proof}
The group $\SL(n,\R)$ acts on $\Symp_1(n,\R)$ by $g \cdot G = gGg^T$;
this preserves $\det G = 1$ since $\det(gGg^T) = (\det g)^2 \det G = 1$.
The action is transitive: given $G \in \Symp_1(n,\R)$, write
$G = \ell \ell^T$ by Cholesky with $\det\ell = 1$ (adjusting the
diagonal entry); then $g = \ell$ satisfies $g \cdot I = G$.
The stabilizer of $I$ is $\{g \in \SL(n,\R) : gg^T = I\} = \SO(n)$.
By the orbit-stabilizer theorem, $\Symp_1(n,\R) \cong \SL(n,\R)/\SO(n)$
as smooth manifolds.
\end{proof}

The space $\SL(n,\R)/\SO(n)$ is a classical Riemannian symmetric space
of noncompact type.
It is complete, simply connected, and has nonpositive sectional
curvature~\cite{Helgason}.

%--------------------------------------------------
\section{Tangent Space and Deformation Directions}
\label{sec:tangent}
%--------------------------------------------------

\begin{proposition}
\label{prop:tangent}
The tangent space of $\calM_n$ at the basepoint $[I]$ is
\[
  T_{[I]}\calM_n \;\cong\; \Symz(n,\R)
  = \{X \in \Sym(n,\R) : \tr X = 0\}.
\]
\end{proposition}

\begin{proof}
A smooth curve in $\calM_n$ through $[I]$ is represented by a smooth
curve $G(t) \in \Symp_1(n,\R)$ with $G(0) = I$.
The constraint $\det G(t) = 1$ gives, upon differentiation at $t=0$,
\[
  0 = \frac{d}{dt}\det G(t)\big|_{t=0}
    = \tr\!\bigl(G(0)^{-1}\dot{G}(0)\bigr)
    = \tr \dot{G}(0).
\]
Since $G(t)$ is symmetric, $\dot{G}(0)$ is symmetric.
Hence $\dot{G}(0) \in \Symz(n,\R)$.
Conversely, for any $X \in \Symz(n,\R)$, the curve
$G(t) = (\det(I + tX))^{-1/n}(I + tX)$ lies in $\Symp_1(n,\R)$
for small $t$ and has $\dot{G}(0) = X$.
\end{proof}

\begin{corollary}
\label{cor:dim}
\[
  \dim \calM_n = \dim \Symz(n,\R) = \frac{n(n+1)}{2} - 1.
\]
For $n = 2, 3, 4, 5$ this gives $2, 5, 9, 14$.
\end{corollary}

\begin{remark}[Correction of a common error]
\label{rem:correction}
It is tempting to argue that the space of scale-free deformation
directions has dimension $n^2-1$, arising from the traceless
endomorphism space $\mathrm{End}_{\tr=0}(V)$.
This argument is incorrect for quadratic forms.
Linear endomorphisms of $V$ deform arbitrary linear maps, including
non-symmetric ones.
Quadratic forms are symmetric bilinear objects; their infinitesimal
deformations must preserve symmetry.
The correct deformation space is $\Symz(n,\R)$, not
$\mathrm{End}_{\tr=0}(V)$.
The dimension $n^2-1$ counts deformations of a linear map
(e.g., a connection or a representation); the dimension
$\frac{n(n+1)}{2}-1$ counts deformations of a quadratic form.
These are distinct geometric settings and must not be conflated.
\end{remark}

%--------------------------------------------------
\section{Cartan Decomposition and Maximal Flats}
\label{sec:cartan}
%--------------------------------------------------

The Lie algebra of $\SL(n,\R)$ is
\[
  \mathfrak{sl}(n,\R) = \{A \in M_n(\R) : \tr A = 0\}.
\]
It admits the \emph{Cartan decomposition}
\[
  \mathfrak{sl}(n,\R) = \mathfrak{k} \oplus \mathfrak{p},
\]
where
\[
  \mathfrak{k} = \mathfrak{so}(n)
    = \{A : A + A^T = 0\},
  \qquad
  \mathfrak{p} = \Symz(n,\R)
    = \{A : A = A^T,\; \tr A = 0\}.
\]
This decomposition satisfies
\[
  [\mathfrak{k},\mathfrak{k}] \subset \mathfrak{k},\quad
  [\mathfrak{k},\mathfrak{p}] \subset \mathfrak{p},\quad
  [\mathfrak{p},\mathfrak{p}] \subset \mathfrak{k},
\]
confirming that $\SL(n,\R)/\SO(n)$ is a symmetric space.
The tangent space $T_{[I]}\calM_n$ is identified with $\mathfrak{p}$.

\begin{definition}[Cartan subalgebra]
A \emph{Cartan subalgebra} of $\mathfrak{p}$ is a maximal abelian
subspace consisting of simultaneously diagonalizable elements.
The standard Cartan subalgebra is
\[
  \mathfrak{a} = \Diagz(n,\R)
  = \{\diag(x_1,\ldots,x_n) : x_i \in \R,\; \textstyle\sum x_i = 0\},
\]
of dimension $n-1$.
\end{definition}

\begin{proposition}
\label{prop:flat}
The exponential image
\[
  F_0 = \exp(\mathfrak{a})
  = \{\diag(e^{x_1},\ldots,e^{x_n}) : \textstyle\sum x_i = 0\}
\]
is a maximal totally geodesic flat submanifold of $\calM_n$,
isometric to $\R^{n-1}$ with the Euclidean metric induced by
$\langle X, Y\rangle = \tr(XY)$.
\end{proposition}

\begin{proof}
Elements of $\mathfrak{a}$ commute; hence the exponential map
restricts to a homomorphism from $(\mathfrak{a}, +)$ into $\calM_n$.
Since $\mathfrak{a}$ consists of simultaneously diagonalizable
real-symmetric matrices (hence with real eigenvalues), the image
$F_0$ is a flat submanifold.
Maximality follows because $\mathfrak{a}$ is a maximal abelian
subspace of $\mathfrak{p}$~\cite[Ch.~V]{Helgason}.
The induced metric on $F_0$ pulls back to $\tr(XY)$ on $\mathfrak{a}$,
which is Euclidean.
\end{proof}

Every maximal flat in $\calM_n$ is of the form $g \cdot F_0$ for some
$g \in \SL(n,\R)$.
The rank of $\calM_n$ is $\dim \mathfrak{a} = n-1$.

%--------------------------------------------------
\section{Sectional Curvature}
\label{sec:curvature}
%--------------------------------------------------

The $\SL(n,\R)$-invariant Riemannian metric on $\calM_n$ is induced
by the Killing form of $\mathfrak{sl}(n,\R)$.
At the basepoint $[I]$, for $X, Y \in \mathfrak{p} = \Symz(n,\R)$,
the inner product is
\[
  \langle X, Y \rangle = \tr(XY).
\]
This is positive definite on $\mathfrak{p}$ (since $\tr(X^2) > 0$
for nonzero symmetric $X$) and $\SO(n)$-invariant.

\begin{proposition}[Sectional curvature]
\label{prop:curvature}
For linearly independent $X, Y \in \mathfrak{p}$, the sectional
curvature of $\calM_n$ in the plane spanned by $X$ and $Y$ is
\[
  K(X,Y) = -\frac{\|[X,Y]\|^2}{\|X\|^2\|Y\|^2 - \langle X,Y\rangle^2},
\]
and norms are taken with respect to $\langle\cdot,\cdot\rangle$.
Since $[X,Y] \in \mathfrak{so}(n)$ is skew-symmetric,
$\tr([X,Y]^2) \leq 0$, so $\|[X,Y]\|^2 := -\tr([X,Y]^2) \geq 0$.
In particular, $K(X,Y) \leq 0$ with equality if and only if
$[X,Y] = 0$.
\end{proposition}

\begin{proof}
The curvature tensor of a symmetric space $G/K$ at the identity coset
is $R(X,Y)Z = -[[X,Y],Z]$ for $X,Y,Z \in \mathfrak{p}$~\cite[Ch.~IV]{Helgason}.
The sectional curvature in the plane $\{X,Y\}$ is
$K(X,Y) = \langle R(X,Y)Y, X\rangle / (\|X\|^2\|Y\|^2 - \langle X,Y\rangle^2)
= \langle [[X,Y],Y], X\rangle / (\cdots)$.
Since $[X,Y] \in \mathfrak{k} = \mathfrak{so}(n)$, one computes
$\langle [[X,Y],Y],X\rangle = -\tr([[X,Y],Y]X) = -\tr([X,Y][Y,X])
= \tr([X,Y]^2) \cdot (-1) \leq 0$,
with equality iff $[X,Y] = 0$.
\end{proof}

\begin{corollary}
$\calM_n \cong \SL(n,\R)/\SO(n)$ is a Hadamard manifold: complete,
simply connected, with sectional curvature $K \leq 0$ everywhere.
Curvature vanishes exactly on maximal flats, where $[X,Y]=0$.
\end{corollary}

%--------------------------------------------------
\section{Geodesics and Weyl Chambers}
\label{sec:weyl}
%--------------------------------------------------

\subsection{Geodesics}

Geodesics through the basepoint $[I]$ in $\calM_n$ are the orbits
of one-parameter subgroups generated by elements of $\mathfrak{p}$:
\[
  \gamma_X(t) = [\exp(tX)], \qquad X \in \Symz(n,\R).
\]
If $X = k D k^{-1}$ with $D = \diag(x_1,\ldots,x_n)$ and
$k \in \SO(n)$, then
\[
  \gamma_X(t) = \bigl[k\,\diag(e^{tx_1},\ldots,e^{tx_n})\,k^{-1}\bigr].
\]
The geodesic distance from $[I]$ to $\gamma_X(t)$ is
$|t|\,\|X\| = |t|\sqrt{\tr(X^2)}$.

\begin{proposition}
The exponential map $\exp_{[I]} : \mathfrak{p} \to \calM_n$,
$X \mapsto [\exp X]$, is a diffeomorphism.
\end{proposition}

\begin{proof}
This is a standard consequence of the Cartan--Hadamard theorem applied
to the nonpositively curved simply connected manifold $\calM_n$.
\end{proof}

\subsection{Weyl chambers}

The Weyl group of $(\mathfrak{sl}(n,\R), \mathfrak{a})$ is
$W \cong S_n$, acting by permuting diagonal entries.

\begin{definition}[Positive Weyl chamber]
\[
  \mathfrak{a}^+ = \{\diag(x_1,\ldots,x_n) \in \mathfrak{a} :
    x_1 \geq x_2 \geq \cdots \geq x_n\}.
\]
\end{definition}

\begin{proposition}
Every point of $\calM_n$ is $\SO(n)$-conjugate to a unique point
in $\exp(\mathfrak{a}^+)$.
The geodesic distance from $[I]$ to $[G]$ equals
$\|\log G_D\|$ where $G_D$ is the unique diagonal representative of
$[G]$ in $\exp(\mathfrak{a}^+)$.
\end{proposition}

\begin{proof}
Every positive definite symmetric matrix is orthogonally diagonalizable
with positive eigenvalues.
Ordering the eigenvalues $e^{x_1} \geq \cdots \geq e^{x_n}$ gives the
unique representative in $\exp(\mathfrak{a}^+)$~\cite[Ch.~V]{Helgason}.
Distance from $[I]$ equals $\sqrt{\sum_i x_i^2} = \|\diag(x_i)\|$.
\end{proof}

Thus the geometry of $\calM_n$ is completely controlled by ordered
logarithmic eigenvalue data.

%--------------------------------------------------
\section{Arithmetic Logarithmic Lattices}
\label{sec:lattice}
%--------------------------------------------------

\subsection{Totally real fields and Dirichlet's theorem}

Let $K$ be a totally real number field of degree $r = [K:\Q]$.
Denote its real embeddings by $\sigma_1,\ldots,\sigma_r : K \hookrightarrow \R$.
The ring of integers $\mathcal{O}_K$ has unit group $\mathcal{O}_K^\times$.
By Dirichlet's unit theorem~\cite{Neukirch},
\[
  \mathcal{O}_K^\times \cong \mu(K) \times \Z^{r-1},
\]
where $\mu(K) = \{\pm 1\}$ for totally real fields with $r \geq 2$.
Fix a system of fundamental units $\varepsilon_1,\ldots,\varepsilon_{r-1}$.

\subsection{Logarithmic embedding}

Define the \emph{logarithmic embedding}
\[
  \Lambda : \mathcal{O}_K^\times \to \R^r, \qquad
  \Lambda(\varepsilon) =
  \bigl(\log|\sigma_1(\varepsilon)|,\ldots,\log|\sigma_r(\varepsilon)|\bigr).
\]
Since $|N_{K/\Q}(\varepsilon)| = \prod_i |\sigma_i(\varepsilon)| = 1$,
the image lies in the hyperplane
\[
  H = \Bigl\{x \in \R^r : \sum_{i=1}^r x_i = 0\Bigr\}.
\]
By Dirichlet's theorem, $\Lambda_K := \Lambda(\mathcal{O}_K^\times)$ is
a lattice of rank $r-1$ in $H$.

\begin{definition}[Regulator]
The \emph{regulator} $R_K$ of $K$ is the covolume of $\Lambda_K$ in
$H$ with respect to the standard Euclidean metric on $\R^r$:
\[
  R_K = \bigl|\det(\Lambda(\varepsilon_j))_{i,j=1}^{r-1}\bigr|,
\]
where the $(i,j)$-entry is $\log|\sigma_i(\varepsilon_j)|$
(omitting one row by the sum-zero constraint).
\end{definition}

\subsection{Embedding into shape space}

Identify $H$ with $\Diagz(r,\R) \subset \mathfrak{a}$ via
\[
  x = (x_1,\ldots,x_r) \;\mapsto\; X = \diag(x_1,\ldots,x_r).
\]
This is an isometry: $\langle X, Y\rangle = \tr(XY) = \sum_i x_i y_i$,
the standard inner product on $H \subset \R^r$.

For each unit $\varepsilon \in \mathcal{O}_K^\times$, define
\[
  A_\varepsilon = \exp\!\bigl(\diag(\Lambda(\varepsilon))\bigr)
  = \diag\!\bigl(|\sigma_1(\varepsilon)|,\ldots,|\sigma_r(\varepsilon)|\bigr)
  \in \Symp_1(r,\R).
\]

\begin{definition}[Arithmetic lattice in shape space]
The \emph{arithmetic lattice} associated with $K$ is
\[
  \Gamma_K = \{[A_\varepsilon] : \varepsilon \in \mathcal{O}_K^\times\}
  \subset \calM_r.
\]
\end{definition}

%--------------------------------------------------
\section{Discrete Flats and Regulator Geometry}
\label{sec:flats}
%--------------------------------------------------

\begin{theorem}[Arithmetic flat embedding]
\label{thm:flatembedding}
$\Gamma_K$ is a discrete subset of the maximal flat
$F_0 = \exp(\Diagz(r,\R)) \subset \calM_r$.
The flat $F_0$ is isometric to $\bigl(\R^{r-1},\,\langle\cdot,\cdot\rangle\bigr)$
where $\langle X,Y\rangle = \tr(XY)$,
and under this isometry the regulator satisfies
\[
  R_K = \mathrm{vol}(\Lambda_K) = \mathrm{covol}(\Gamma_K \text{ in } F_0).
\]
\end{theorem}

\begin{proof}
By construction, $[A_\varepsilon] = [\exp(\diag(\Lambda(\varepsilon)))]
\in F_0$ for every $\varepsilon$.
The map $\varepsilon \mapsto [A_\varepsilon]$ is the composition of
$\Lambda$ (injective modulo torsion) with the exponential map
(a diffeomorphism on $\Diagz(r,\R)$); hence $\Gamma_K$ is discrete.
The isometry $F_0 \cong \R^{r-1}$ preserves the inner product
$\tr(XY)$, which restricts to the standard Euclidean inner product on
$\Diagz(r,\R) \cong H$.
Covolume is preserved under isometries, so
$\mathrm{covol}(\Gamma_K) = \mathrm{covol}(\Lambda_K) = R_K$.
\end{proof}

\begin{theorem}[Flat rigidity]
\label{thm:rigidity}
Let $F_1$ and $F_2$ be maximal totally geodesic flats in $\calM_r$.
If $F_1 \cap F_2$ contains an open set, then $F_1 = F_2$.
Consequently, the arithmetic flat $F_0$ is uniquely determined by the
tangent Cartan subalgebra $\mathfrak{a} = \Diagz(r,\R)$.
\end{theorem}

\begin{proof}
Maximal flats correspond to maximal abelian subspaces
of $\mathfrak{p} = \Symz(r,\R)$
under the $\SO(r)$-action.
If $F_1$ and $F_2$ share an open subset,
their tangent spaces coincide at some interior point.
Two maximal abelian subspaces of $\mathfrak{p}$ that agree on
an open set in their exponential images must coincide:
any element of a maximal abelian subspace $\mathfrak{a}$
is determined by its action on a neighborhood of the identity
via the exponential map, and a maximal abelian subspace containing
a given regular element is unique~\cite[Ch.~V, Thm.~6.2]{Helgason}.
Hence the tangent subspaces of $F_1$ and $F_2$ agree globally,
so $F_1 = F_2$.
\end{proof}

\begin{remark}
Theorem~\ref{thm:rigidity} ensures that distinct totally real fields
of the same degree, if non-isomorphic, give rise to distinct arithmetic
flats inside $\calM_r$ (up to the $\SL(r,\R)$-action).
This makes the geometric data of the flat — its shape, covolume, and
position — an arithmetic invariant of the field.
\end{remark}

%--------------------------------------------------
\section{Explicit Examples}
\label{sec:examples}
%--------------------------------------------------

\subsection{Real quadratic fields ($r = 2$)}

For $r = 2$, the symmetric space is $\calM_2 \cong \SL(2,\R)/\SO(2)$,
the hyperbolic plane.
The Cartan subalgebra $\mathfrak{a} = \Diagz(2,\R)$ is one-dimensional,
spanned by $\diag(1,-1)$.
The corresponding flat $F_0$ is a single geodesic.

Let $K = \Q(\sqrt{d})$ with $d > 0$ squarefree, and let
$\varepsilon > 1$ be the fundamental unit.
The logarithmic embedding gives
$\Lambda(\varepsilon) = (\log\varepsilon, -\log\varepsilon) \in H \cong \R$.
The arithmetic lattice $\Gamma_K \subset F_0$ is generated by
\[
  [A_\varepsilon] = \bigl[\diag(\varepsilon, \varepsilon^{-1})\bigr],
\]
at distance $\sqrt{2}\log\varepsilon$ from $[I]$ along the geodesic.
The regulator is $R_K = \log\varepsilon$.

\subsection{The golden ratio ($K = \Q(\sqrt{5})$)}

The field $\Q(\sqrt{5})$ has discriminant $5$ and fundamental unit
$\varphi = \frac{1+\sqrt{5}}{2}$, satisfying $\varphi^2 = \varphi + 1$.
Its conjugate is $\hat\varphi = \frac{1-\sqrt{5}}{2}$, with
$\varphi\hat\varphi = -1$, so $|\hat\varphi| = \varphi^{-1}$.

The logarithmic embedding gives
$\Lambda(\varphi) = (\log\varphi, -\log\varphi)$,
and the arithmetic lattice is generated by
\[
  [A_\varphi] = \bigl[\diag(\varphi, \varphi^{-1})\bigr]
  \in F_0 \subset \calM_2.
\]
The translation length is $\sqrt{2}\log\varphi \approx 0.6624$,
and the regulator is $R_{\Q(\sqrt{5})} = \log\varphi$.

Thus the golden ratio $\varphi$ appears as the generator of a geodesic
translation in the hyperbolic plane $\SL(2,\R)/\SO(2)$.
This is not a coincidence of numerology but a structural consequence:
$\varphi$ is the fundamental unit of $\Q(\sqrt{5})$, and the
hyperbolic plane is the shape space of rank-$2$ quadratic forms.

\subsection{Totally real cubic fields ($r = 3$)}

For $r = 3$, the symmetric space is $\calM_3 \cong \SL(3,\R)/\SO(3)$
of dimension $5$ and rank $2$.
A totally real cubic field $K$ has two fundamental units
$\varepsilon_1, \varepsilon_2$, giving logarithmic vectors
$\mathbf{x}_j = \Lambda(\varepsilon_j) \in H \subset \R^3$ for $j=1,2$.
These generate a rank-$2$ lattice in the two-dimensional flat $F_0$,
with regulator
\[
  R_K = \bigl|\det\begin{pmatrix}
    \log|\sigma_1(\varepsilon_1)| & \log|\sigma_1(\varepsilon_2)| \\
    \log|\sigma_2(\varepsilon_1)| & \log|\sigma_2(\varepsilon_2)|
  \end{pmatrix}\bigr|.
\]
The arithmetic lattice $\Gamma_K$ is a two-dimensional discrete set
inside the flat $F_0 \cong \R^2$, with covolume $R_K$.
The shape (aspect ratio) of the lattice encodes arithmetic properties
of $K$ beyond the regulator itself.

%--------------------------------------------------
\section{Conclusion}
\label{sec:conclusion}
%--------------------------------------------------

We have established the following results for scale-free quadratic
moduli space $\calM_n = \Symp(n,\R)/\R_{>0}$:

\begin{enumerate}
\item \emph{Symmetric space identification:}
      $\calM_n \cong \SL(n,\R)/\SO(n)$, proved via the transitive
      $\SL(n,\R)$-action on the determinant-one slice and the
      orbit-stabilizer theorem.

\item \emph{Correct tangent space:} $T_{[I]}\calM_n \cong \Symz(n,\R)$,
      of dimension $\frac{n(n+1)}{2}-1$.
      The incorrect dimension $n^2-1$, arising from traceless linear
      endomorphisms, is explicitly excluded: quadratic forms admit
      only symmetric deformations.

\item \emph{Riemannian geometry:} The trace metric $\langle X,Y\rangle
      = \tr(XY)$ makes $\calM_n$ a Hadamard manifold with sectional
      curvature $K(X,Y) = -\|[X,Y]\|^2/\|X\wedge Y\|^2 \leq 0$.
      Geodesics are one-parameter subgroups; Weyl chambers encode
      eigenvalue ordering.

\item \emph{Arithmetic flat embedding:} For a totally real field $K$
      of degree $r$, the logarithmic unit lattice embeds as a discrete
      subset $\Gamma_K$ of the maximal flat $F_0 \subset \calM_r$,
      with regulator $R_K$ equal to the Euclidean covolume.

\item \emph{Flat rigidity:} The arithmetic flat $F_0$ is uniquely
      determined by its tangent Cartan subalgebra $\Diagz(r,\R)$.

\item \emph{Golden ratio example:} The fundamental unit $\varphi$ of
      $\Q(\sqrt{5})$ generates a geodesic translation of length
      $\sqrt{2}\log\varphi$ in $\calM_2 \cong \SL(2,\R)/\SO(2)$.
\end{enumerate}

\bibliographystyle{plain}
\begin{thebibliography}{99}

\bibitem{Helgason}
S.~Helgason,
\emph{Differential Geometry, Lie Groups, and Symmetric Spaces},
Academic Press, 1978.

\bibitem{Neukirch}
J.~Neukirch,
\emph{Algebraic Number Theory},
Springer, 1999.

\bibitem{BridsonHaefliger}
M.~R.~Bridson and A.~Haefliger,
\emph{Metric Spaces of Non-Positive Curvature},
Springer, 1999.

\bibitem{Ratcliffe}
J.~G.~Ratcliffe,
\emph{Foundations of Hyperbolic Manifolds},
Springer, 2006.

\bibitem{Eberlein}
P.~B.~Eberlein,
\emph{Geometry of Nonpositively Curved Manifolds},
University of Chicago Press, 1996.


\bibitem{Gentry:CKM}
M.~L. Gentry,
\newblock {CKM} quark mixing from geodesic axes of hyperbolic 3-manifold holonomy,
\newblock \emph{Phys. Rev. D} (2026), submitted March 2026, temporary {ID} es2026mar11\_966.

\bibitem{Gentry:PMNS}
M.~L. Gentry,
\newblock Lepton mixing from Borel structure of hyperbolic holonomy,
\newblock \emph{Phys. Rev. D} (2026), submitted March 2026, temporary {ID} es2026mar13\_942.
\end{thebibliography}

\end{document}
